
\begin{ledgroupsized}[r]{120mm}
\footnotesize
\pstart
\noindent\textbf{Überlieferung:}
\pend
\end{ledgroupsized}
\begin{ledgroupsized}[r]{114mm}
\footnotesize
\pstart
\parindent -6mm
\makebox[6mm][l]{\textit{L}}%
Aufzeichnung: LH XXXVII 5 Bl. 108. Ein Blatt 8\textsuperscript{o}; Papiererhaltungsmaßnahmen. Zwei Seiten.
\pend
\end{ledgroupsized}
\selectlanguage{latin}		%
\vspace{8mm}
\pstart
\noindent
\href{\Scan{LH_037_05_108r}}%
{\lbrack108~r\textsuperscript{o}\rbrack}%
\pend
\begin{edfigure}[1.5em]%
\vspace{-1em}
\hfill%
\hspace*{5mm}%
\includegraphics[scale=1]{%
gesamttex/edit_VIII-4/images/LH_037_05_108_d1_108r.pdf}%
\hfill%
\includegraphics[scale=1]{%
gesamttex/edit_VIII-4/images/LH_037_05_108_d2_108r.pdf}%
\hfill%
\edcaption[1em]{%
\lbrack\textit{Fig.~1}\rbrack%
\hspace*{41mm}%
\lbrack\textit{Fig.2}\rbrack%
\hspace*{5mm}
\edtext{}{%
\lemma{}%
\Afootnote[nonum]{%
\edlabel{LH_37_05_108_05a}%
\textit{Am Rand von \bracks{Fig.~2}:}\\
\textit{PF} vel \textit{P\greek{f}} %
$=N-\varepsilon$ vel $\varepsilon-N$\bracks{,} %
$D=Q\langle D\rangle$\bracks{,} %
sed $M=GF$\bracks{,} %
$GF-QP= %
FP+GQ= %
FP+CT=\text{bis}\;CT$.
Ergo $M-D=\text{bis}\;\overline{N-\varepsilon}$.%
\edlabel{LH_37_05_108_05b}%
}}%
}%
\end{edfigure}%
\pstart%
\noindent%
\edtext{%
Capacitas totius \textit{A} partis \textit{PQ}
\edtext{sit \textit{B} spatium}{%
\lemma{sit \textit{B}}%
\Bfootnote{%
\textit{(1)}~quantitas
\textit{(2)}~spatium~\textit{L}}}
hydrargyri\indexS{spatium hydrargyri} \textit{G}
\edtext{et \textit{QX} sit \textit{H}}{%
\lemma{}%
\Bfootnote{et \textit{QX} sit \textit{H} \textit{erg.~L}}}%
\lbrack ,\rbrack\ \textit{PX} sit $B-H$ %
\edtext{et}{\lemma{}\Bfootnote{et \textit{erg.~L}}}
erit spatium\indexS{spatium aeris} aeris
\edtext{primum \textit{STV},}{%
\lemma{}%
\Bfootnote{%
\bars{primum \textit{erg.}}
\textit{(1)}~\textit{AC}
\textit{(2)}~\textit{SV}
\textit{(3)}~\textit{STV}~\textit{L}}}%
erit $A-G-B+H$. 
}{%
\lemma{Capacitas \bracks{...} $A-G-B+H$.}%
\Cfootnote{Die Buchstaben beziehen sich womöglich auf teils gestrichene oder ersetzte Bezeichnungen in einer früheren Fassung von \bracks{\textit{Fig.~2}}.}
}
Sed spatium %
\edtext{aeris\indexS{spatium aeris} mutatione}{%
\lemma{aeris}%
\Bfootnote{%
\textit{(1)}~post m
\textit{(2)}~mutatione~\textit{L}}}
facta
\edtext{est $A-G-B+C$. Et}{%
\lemma{est $A-G-B+C$.}%
\Bfootnote{%
\textit{(1)}~Sunt ergo
\textit{(2)}~Ergo
\textit{(3)}~Et~\textit{L}}}
densitas\indexS{densitas aeris} mutatione
\edtext{facta est}{%
\lemma{}%
\Bfootnote{%
facta \bars{\textit{AG} \textit{gestr.}} 
est~\textit{L}}}
ad priorem ut spatium prius aeris\indexS{spatium aeris} ad posterius. Sed pondus
\begin{tabular}{c}
posterius\\ prius
\end{tabular}
est ad
\begin{tabular}{c}
prius\\ posterius
\end{tabular}
\edtext{ut \textit{C} ad \textit{H}. }{%
\lemma{ut}%
\Bfootnote{%
\textit{(1)}~\textit{C} ad \textit{I}
\textit{(2)}~\textit{H} ad \textit{C}
\textit{(3)}~\textit{C} ad \textit{H}~\textit{L}}}%
Et calores\indexS{calor} sunt in ratione directa ponderum\indexS{ratio ponderum} et reciproca
\edtext{densitatum\indexS{densitas} seu ob spatia}{%
\lemma{densitatum}%
\Bfootnote{%
\textit{(1)}~. Ergo calor
\textit{(a)}~pri
\textit{(b)}~\bars{posterior est \textit{streicht Hrsg.}}
 ad priorem
\textit{(2)}~seu ob spatia~\textit{L}}}
hic aeris densitatibus\indexS{densitas aeris} reciproce proportionalia, sunt in ratione ponderum\indexS{ratio ponderum} et spatiorum\indexS{ratio spatiorum}. Ergo %
\edtext{calor posterior est ad priorem, ut}{%
\lemma{spatiorum\indexS{ratio spatiorum}.}%
\Bfootnote{%
\textit{(1)}~prior est ad posteriorem, ut
\textit{(2)}~posterior \bracks{...}, ut~\textit{L}}}	
\indexS{calor} 
\textit{H} in $\overline{A-G-B+H}$ est ad \textit{C} in $\overline{A-G-B+C}$.
\pend
\pstart
\edlabel{LH_37_05_108_04a}%
\rule[0mm]{0pt}{6mm}
\edtext{%
$QXVT=A$, 100, $FQ=L=60$. 
\edtext{%
$TC=20=N$.%
}{%
\lemma{}%
\Afootnote{%
\textit{Am Rand: $\displaystyle\efrac{130}{65}$}
}}
$B=100\PM 20$\lbrack ,\rbrack\ $M=60%
\edtext{\PM \text{bis}\;20$\bracks{,} \textit{BM}}{%
\lemma{}%
\Bfootnote{%
$\PM \text{bis}\, 20$\bars{$\displaystyle\frac{30}{100}.$ \textit{streicht Hrsg.}} 
\textit{BM}~\textit{L}}}
$=100\cdot 60+\text{bis}\, 20^2\PM%
\begin{array}{l}
20\cdot 60\\
\text{bis}\, 20\cdot 100
\end{array}%
\edtext{%
\!=100\cdot 60\cdot 2$\;%
}{%
\lemma{}%
\Bfootnote{%
$=100\cdot 60\cdot 2$ \bars{seu $\displaystyle20^2\,\PM\,\frac{60}{2}$ \textit{erg. u. gestr.}} 
et~\textit{L}}}
et $20^2\,
\edtext{\PM\;\langle60:\text{duo}\cdot 100\cdot 20\rangle%
}{%
\lemma{$\PM\;\langle60:\text{duo}\cdot 100\cdot 20\rangle$}%
\Cfootnote{Der Term müsste 2600 ergeben, damit die Gleichung erfüllt ist.}}
=100\cdot 60\; \overline{2-1}:\text{duo}$. Sit $60:\text{duo}+100 = \text{bis}\, 65$\bracks{,} %
\edtext{%
Calor\indexS{calor} $C=2$\bracks{,} %
}{\lemma{}\Bfootnote{Calor $C=2$ \textit{erg.~L}}}
$65=V$ et $1:2=\varepsilon$
 et $\displaystyle\frac{2-1}{2}=\PM\, \overline{1:2}$, et fit $20^2\PM \text{bis}\, 65\cdot 20=\PM\, \overline{1:2}\cdot 100\cdot 60$ et $20=\surd\llcorner 65^2\PM\, \overline{1:2}\cdot 100\cdot 60\lrcorner \MP 65$
}{%
\lemma{}%
\Afootnote{%
\textit{Nebenrechnungen am Rand von \bracks{Fig.~1} und \bracks{Fig.~2}:}\\%
\rule[0mm]{0pt}{6mm}%
$\begin{array}[t]{r}
120\phantom{00}\\00\\\overline{12000}
\end{array}$
$\begin{array}[t]{r}
130\phantom{0}\\20\\\overline{2600}\\400\\\overline{3000}
\end{array}$
$\begin{array}[t]{r}
400\\2600\\\overline{3000}
\end{array}$
$\begin{array}[t]{r}
65\\65\\\overline{325}\\390\phantom{0}\\\overline{4225}
\end{array}$
$\begin{array}[t]{r}
4225\\3000\\7225\\\underline{\phantom{7}\ldotp\phantom{2}\ldotp}\\\cancel{825}\\\cancel{6}\phantom{5}\\\cancel{1}\phantom{85}
\end{array}$
$\begin{array}[t]{r}
6000\\100\\200\\4000\\\overline{10300}
\end{array}$
$\begin{array}[t]{r}
6000\\800\\1200\\4000\\\overline{12000}
\end{array}$%
}}%
\edlabel{LH_37_05_108_04b}
\pend
\vspace{0.5em}
\pstart
\href{\Scan{LH_037_05_108v}}%
{\lbrack108~v\textsuperscript{o}\rbrack} 
$\displaystyle\frac{\text{cal. nov.}}{\text{cal. vet.}}\cdot AL=AL\PM AN+2NN$. %
\edtext{Ut ergo}{%
\lemma{$=AL\PM 2AN+2NN$.}%
\Bfootnote{%
\textit{(1)}~\bars{Fiet quib \textit{streicht Hrsg.}}
\textit{(2)}~$NN\PM\, \displaystyle\frac{1}{2}AN+cal.$
\textit{(3)}~Si ergo
\textit{(4)}~Ut ergo~\textit{L}}}
gradus
\edtext{ascribamus, ponamus}{%
\lemma{ascribamus,}%
\Bfootnote{%
\textit{(1)}~fiet
\textit{(2)}~ponamus~\textit{L}}}
\rule[0mm]{0pt}{4mm}%
datum numerum partium quibus distat \textit{F} a \textit{P}. 
\pend%
\pstart%
\noindent%
\rule[0mm]{0pt}{6mm}%
$\text{cal. nov.}\edlabel{LH_37_05_108v_1i}=\text{cal. vet.}\;\PM\;\text{cal. vet.}\;A:L\;\PM$%
\edtext{}{%
\xxref{LH_37_05_108v_1i}{LH_37_05_108v_1f}
\lemma{$=\text{cal. vet.}\;\PM\;\text{cal. vet.}\;A:L \PM$}%
\Bfootnote{%
\textit{(1)}~$L+N$
\textit{(2)}~$\displaystyle\frac{\text{cal. nov.}-\text{cal. vet.}}{2}+$
\textit{(3)}~$\displaystyle\frac{\text{cal. nov.}-\text{cal. vet.}}{2}$ sit $\PM\;$
\textit{(a)}~\textit{D}
\textit{(b)}~$\varepsilon$~\textit{L}}}%
$\displaystyle\frac{\text{cal. nov.}-\text{cal. vet.}}{2}$ sit $\PM\;\varepsilon\,$\edlabel{LH_37_05_108v_1f} %
et $\displaystyle\frac{A+2L}{2}$ sit 2\textit{Y}. %
\rule[0mm]{0pt}{6mm}%
Fiet $\displaystyle NN\;\PM\; 2YN+YY=YY+\varepsilon\cdot AL$ et $N\PM Y=\sqrt{YY\PM \varepsilon\cdot AL}$ seu $\displaystyle\frac{2A+L}{4}=Y$\lbrack ,\rbrack\ $N=\sqrt{YY\PM \varepsilon\cdot AL}\MP Y$.%
\rule[0mm]{0pt}{6mm}%
\pend
\pstart
\rule[0mm]{0pt}{6mm}%
\edtext{\textit{A} volumen aeris\indexS{volumen aeris} ordinarii, \textit{B} alterati, \textit{L} volumen \mercury\indexS{volumen mercurii} ordinarii, \textit{M} alterati. \edlabel{LH_37_05_108_01a}%
\edtext{In discursu absoluto}{%
\lemma{In discursu absoluto}%
\Cfootnote{%
Im gültigen Text von N.~\ref{60270} bezeichnet Leibniz diese Größe mit \textit{V} (S.~\refpassageohne{LH_37_04_169-170_17a}{LH_37_04_169-170_17b} u.ö.), in den gestrichenen Textteilen mit \textit{Y} (S.~\refpassageohne{LH_37_04_169-170_16a}{LH_37_04_169-170_16b} u.ö.).}}
pro \textit{Y} scripsi \textit{V}\edlabel{LH_37_05_108_01b}.}{%
\lemma{}\Bfootnote{%
\textit{A} volumen aeris\indexS{volumen aeris} \bracks{...} absoluto pro \textit{Y}
\bars{scripsi \textit{streicht Hrsg.}}
scripsi \textit{V}. \textit{erg.~L}}}
\pend
\pstart
\noindent
\rule[0mm]{0pt}{6mm}%
$BM:AL\squaredots\, C:1$.%
\qquad%
$\PM\:B\,\MP\,A=N$%
\qquad%
$\PM\:M\,\MP\,L=2N$%
\qquad%
$B=\PM\:N+A$
\qquad%
$M=\PM\:2N+L$\\%
\rule[0mm]{0pt}{6mm}%
\edlabel{LH_37_05_108_02a}%
$BM=AL\,\PM\,2AN+
\edtext{LN}{\lemma{}\Bfootnote{\textit{LN} \textit{erg.~L}}}
+2NN=ALC$ seu $NN\,\PM\,\displaystyle\frac{2A+L}{2}\,N+\displaystyle\frac{1}{2}\,AL=\displaystyle\frac{1}{2}\,ALC$ seu %
\rule[0mm]{0pt}{6mm}%
 $NN\,\PM \displaystyle\frac{2A+L}{2}\,N=\displaystyle\frac{C-1}{2}\,AL$. Sit $\displaystyle\frac{2A+L}{4}=Y$ et dimidia diff. inter \textit{C} et 1 sit $\varepsilon$.%
\edlabel{LH_37_05_108_02b}
\\%
\rule[0mm]{0pt}{7mm}%
Patet in casu $\PM =+$ esse $\displaystyle\frac{C-1}{2}=+\:\varepsilon$\lbrack ,\rbrack\ in casu $\MP =+$ esse $\displaystyle\frac{C-1}{2}=-\:\varepsilon$. Ergo $\displaystyle\frac{C-1}{2}=\PM\:\varepsilon$ et fit\\%
\rule[0mm]{0pt}{6mm}%
$NN\,\PM\;2\,YN=\PM\;\varepsilon\,AL$ et $N=(\PM )\sqrt{YY\PM\: \varepsilon AL}\;\MP\,Y$ vel %
\edlabel{LH_37_05_108_03a}%
$\PM\;N=(\PM)\sqrt{\dots}-Y$\lbrack ,\rbrack\ $N=(\PM)\:\PM \sqrt{\dots}\:\MP\;Y$\lbrack ,\rbrack\ %
\rule[0mm]{0pt}{5mm}%
id est $N=\PM \sqrt{\dots}\MP\,Y$. 
\edlabel{LH_37_05_108_03b}%
\edtext{}{%
\xxref{LH_37_05_108_03a}{LH_37_05_108_03b}%
\lemma{}%
\Cfootnote{%
Auslassungspunkte von Leibniz.}}%
\edtext{Nam}{%
\lemma{}%
\Bfootnote{%
$=\PM \sqrt{\dots}\MP Y$.
\bars{Pro \textit{Y} in discursu absoluto scripsi \textit{V}. \textit{gestr.}} 
Nam~\textit{L}}}
in casu $\varepsilon=0$, oportet et \textit{N} fieri 0, unde tunc si %
\rule[0mm]{0pt}{5mm}%
$N\,\MP =+$ debet esse $(\PM)=-$, et contra.%
\pend
\vspace{1.5em}
\pstart
\noindent
\bracks{\textit{Nachfolgend kleingedruckter Text in L quer zur Schreibrichtung und gestrichen:\,}}
\pend
\nobreak
\vspace{0.5em}
\footnotesize	%
\pstart 
\noindent
\rule[0mm]{0pt}{5.5mm}
\hspace{4mm}$\displaystyle\frac{A}{B}=\displaystyle\frac{C}{D}$\quad$\displaystyle\frac{L}{M}=\displaystyle\frac{A}{B}$ si $C=D$\\%
\rule[0mm]{0pt}{5.5mm}
et $\displaystyle\frac{L}{M}=\displaystyle\frac{C}{D}$ si $A=B$\\%
\rule[0mm]{0pt}{5.5mm}
et $\displaystyle\frac{(L)}{(M)}=\displaystyle\frac{(C)}{(D)}$ si $(A)=(B)$.\\%
\rule[0mm]{0pt}{5.5mm}
Ergo $\displaystyle\frac{((L))}{((B))}=\efrac{((A))}{}$ %
\quad%
\normalsize%
\bracks{\textit{Text bricht ab.}}
\pend

